\section{Workspace Layout and Dependencies}

The environment was prepared in a single working directory, separating the main frameworks from the PLEXE subprojects dedicated to VLC and 5G. This choice facilitates both incremental builds and traceability of the versions used.

\begin{lstlisting}[caption={Reference layout of the local workspace},label={lst:workspace_layout}]
~/src/
  omnetpp-6.2/
  sumo/
  inet/
  simu5g/
  veins/
    subprojects/
      veins_vlc/
      veins_inet
  plexe/
    subprojects/
      plexe_vlc/
      plexe_5g/
      plexe_hetnet/
\end{lstlisting}

Basic dependencies and compatibilities were determined by cross-checking the Plexe \emph{Building} page, official tutorials, the README files of the \texttt{plexe\_vlc} and \texttt{plexe\_5g} subprojects, the SUMO repository, and the OMNeT++ download page \cite{plexe_building,plexe_tutorial,plexe_vlc_readme,plexe_5g_readme,sumo_repo,omnet_download}. In particular:
\begin{itemize}
\item for SUMO, the adopted workflow was latest-version dependencies plus checkout of version 1.20 before compilation;
\item for OMNeT++ 6.2, the \texttt{./install.sh} script from the source package was used, which automates system dependencies, Python packages, and virtual environment setup \cite{omnet_download};
\item for Veins, Plexe, Simu5G, INET, and the Plexe subprojects, the official operational instructions from the respective README files and tutorials were followed \cite{plexe_building,plexe_vlc_readme,plexe_5g_readme}.
\end{itemize}
